\emergencystretch=3em
\renewcommand{\vec}[1]{#1}
\renewcommand{\solution}{\noindent Solution: }
\begin{enumerate}[label=\thesubsection.\arabic*, ref=\thesubsection.\theenumi]
\item For real $a$, the principal values satisfy
\begin{align}
	\tan^{-1} a + \cot^{-1} a = \frac{\pi}{2} \label{eq:EE26BTECH11205-tan-cot}
\end{align}
\item If $\vec{A}$ is a square matrix, $\vec{I}$ commutes with $\vec{A}$, so
\begin{align}
	\brak{\vec{I}+\vec{A}}^2 = \vec{I} + 2\vec{A} + \vec{A}^2 \label{eq:EE26BTECH11205-idem-binom}
\end{align}
\item A system of $n$ linear equations in $n$ unknowns is written as
\begin{align}
	\vec{A}\vec{x} = \vec{b} \label{eq:EE26BTECH11205-linsys}
\end{align}
It has a unique solution $\vec{x} = \vec{A}^{-1}\vec{b}$ if and only if $\mydet{\vec{A}} \neq 0$. The solution is obtained by reducing the augmented matrix $\myvec{\vec{A} & \vec{b}}$ to an upper triangular form using row operations (Gaussian elimination) followed by back substitution.
\item Determinants.
\begin{enumerate}
\item Adding a multiple of one row to another row does not change the determinant.
\begin{align}
	\mydet{\vec{E}\vec{M}} = \mydet{\vec{M}},\quad \vec{E} = \vec{I} + \lambda \vec{e}_i\vec{e}_j^{\top},\; i \neq j \label{eq:EE26BTECH11205-det-elem}
\end{align}
\solution $\vec{E}$ is triangular with unit diagonal, so $\mydet{\vec{E}} = 1$ and $\mydet{\vec{E}\vec{M}} = \mydet{\vec{E}}\mydet{\vec{M}}$.
\item The determinant of a triangular matrix is the product of its diagonal entries.
\begin{align}
	\mydet{\vec{U}} = \prod_{i=1}^{n} u_{ii} \label{eq:EE26BTECH11205-det-tri}
\end{align}
\item If a row of $\vec{M}$ is a multiple of another row, then the row operation in \eqref{eq:EE26BTECH11205-det-elem} produces a zero row, and
\begin{align}
	\mydet{\vec{M}} = 0 \label{eq:EE26BTECH11205-det-prop}
\end{align}
\end{enumerate}
\item For differentiable $f$,
\begin{align}
	\int e^x \brak{f(x) + f^{\prime}(x)}dx = e^x f(x) + c \label{eq:EE26BTECH11205-exp-int}
\end{align}
\solution By the product rule, $\frac{d}{dx}\brak{e^x f(x)} = e^x\brak{f(x) + f^{\prime}(x)}$.
\item The area of the parallelogram with adjacent sides $\vec{a}$ and $\vec{b}$ is
\begin{align}
	\sqrt{\det\brak{\vec{M}^{\top}\vec{M}}}, \quad \vec{M} = \myvec{\vec{a} & \vec{b}} \label{eq:EE26BTECH11205-area-gram}
\end{align}
\item The direction cosines of a line with direction vector $\vec{d}$ are the components of the unit vector
\begin{align}
	\frac{\vec{d}}{\norm{\vec{d}}} = \myvec{\cos\alpha \\ \cos\beta \\ \cos\gamma} \label{eq:EE26BTECH11205-dircos}
\end{align}
\item A linear fractional (M\"obius) function and its matrix are related by
\begin{align}
	f(x) = \frac{ax+b}{cx+d} \longleftrightarrow \vec{M} = \myvec{a & b \\ c & d} \label{eq:EE26BTECH11205-mobius}
\end{align}
Then $f_{\vec{M}}\circ f_{\vec{N}} = f_{\vec{M}\vec{N}}$ and $f_{k\vec{M}} = f_{\vec{M}}$ for $k \neq 0$ wherever the functions are defined.

\smallskip
\solution In homogeneous coordinates,
\begin{align}
	\vec{M}\myvec{x \\ 1} = \myvec{ax+b \\ cx+d} = \brak{cx+d}\myvec{f(x) \\ 1}
\end{align}
so applying $f_{\vec{N}}$ followed by $f_{\vec{M}}$ multiplies the column vector by $\vec{N}$ and then by $\vec{M}$, up to non-zero scalars, which do not change the ratio of the entries.
\item If $f\circ f = \text{identity}$ on its domain and $f$ maps the domain onto itself, then $f$ is one-one, onto, and
\begin{align}
	f^{-1} = f \label{eq:EE26BTECH11205-involution}
\end{align}
\item For angles $\alpha, \beta$, multiplying unit complex numbers adds their angles:
\begin{align}
	\brak{\cos\alpha + \mathrm{j}\sin\alpha}\brak{\cos\beta + \mathrm{j}\sin\beta} = \cos\brak{\alpha+\beta} + \mathrm{j}\sin\brak{\alpha+\beta} \label{eq:EE26BTECH11205-euler-mult}
\end{align}
\item For $uv < 1$,
\begin{align}
	\tan^{-1}u + \tan^{-1}v &= \tan^{-1}\frac{u+v}{1-uv} \label{eq:EE26BTECH11205-arctan-add}
	\\
	\implies 2\tan^{-1}u &= \tan^{-1}\frac{2u}{1-u^2}, \quad \abs{u} < 1 \label{eq:EE26BTECH11205-arctan-double}
\end{align}
\solution For $\alpha = \tan^{-1}u$, $\beta = \tan^{-1}v$, $\tan\brak{\alpha+\beta} = \frac{u+v}{1-uv}$. If $uv<1$, then $\alpha+\beta \in \brak{-\frac{\pi}{2},\frac{\pi}{2}}$, which is the range of $\tan^{-1}$.
\item For real $u(x) \neq 0$,
\begin{align}
	\frac{d}{dx}\ln\abs{u} = \frac{u^{\prime}}{u} \label{eq:EE26BTECH11205-logdiff}
\end{align}
\item
\begin{align}
	\frac{d}{dx}\cos^{-1}x = -\frac{1}{\sqrt{1-x^2}}, \quad -1<x<1 \label{eq:EE26BTECH11205-dacos}
\end{align}
\item If $f^{\prime}(x) \ge 0$ on an interval and $f^{\prime}(x) = 0$ only at isolated points, then $f$ is strictly increasing on that interval. \label{eq:EE26BTECH11205-monotone}
\item The gradient of $F(x,y)$ at a point $\vec{P}$ on the curve $F(x,y) = 0$ is normal to the curve. The normal line has direction $\nabla F\brak{\vec{P}}$:
\begin{align}
	\nabla F = \myvec{\frac{\partial F}{\partial x} \\ \frac{\partial F}{\partial y}}, \quad \vec{x} = \vec{P} + \lambda\nabla F\brak{\vec{P}} \label{eq:EE26BTECH11205-grad-normal}
\end{align}
\item Integration by parts:
\begin{align}
	\int u\,dv = uv - \int v\,du \label{eq:EE26BTECH11205-ibp}
\end{align}
\item The linear differential equation $\frac{dy}{dx} + P(x)y = Q(x)$ has the integrating factor $\mu = e^{\int P\,dx}$ and
\begin{align}
	\frac{d}{dx}\brak{\mu y} = \mu Q \label{eq:EE26BTECH11205-linode}
\end{align}
\solution $\frac{d}{dx}\brak{\mu y} = \mu\brak{y^{\prime} + Py}$ because $\mu^{\prime} = P\mu$.
\item A separable equation $f(x)\,dx + g(y)\,dy = 0$ is solved by
\begin{align}
	\int f(x)\,dx + \int g(y)\,dy = C \label{eq:EE26BTECH11205-sepode}
\end{align}
\item The vector $\vec{\beta}$ decomposes as $\vec{\beta} = \vec{\beta}_1 + \vec{\beta}_2$ with $\vec{\beta}_1 \parallel \vec{\alpha}$ and $\vec{\beta}_2 \perp \vec{\alpha}$, where
\begin{align}
	\vec{\beta}_1 = \vec{P}\vec{\beta},\quad \vec{\beta}_2 = \brak{\vec{I}-\vec{P}}\vec{\beta},\quad \vec{P} = \frac{\vec{\alpha}\vec{\alpha}^{\top}}{\vec{\alpha}^{\top}\vec{\alpha}} \label{eq:EE26BTECH11205-proj}
\end{align}
\solution Let $\vec{\beta}_1 = \lambda\vec{\alpha}$. Then $\vec{\alpha}^{\top}\vec{\beta}_2 = \vec{\alpha}^{\top}\brak{\vec{\beta} - \lambda\vec{\alpha}} = 0$ gives $\lambda = \frac{\vec{\alpha}^{\top}\vec{\beta}}{\vec{\alpha}^{\top}\vec{\alpha}}$.
\item A line with direction $\vec{d}$ is parallel to the planes $\vec{n}_1^{\top}\vec{x} = c_1$ and $\vec{n}_2^{\top}\vec{x} = c_2$ if
\begin{align}
	\myvec{\vec{n}_1^{\top} \\ \vec{n}_2^{\top}}\vec{d} = \vec{0} \implies \vec{d} = \vec{n}_1\times\vec{n}_2 \label{eq:EE26BTECH11205-line-planes}
\end{align}
The line through $\vec{P}$ is $\vec{x} = \vec{P} + \lambda\vec{d}$, and its Cartesian form is $\frac{x-p_1}{d_1} = \frac{y-p_2}{d_2} = \frac{z-p_3}{d_3}$.
\item If $X$ is the number of successes in $n$ independent trials, each with probability $p$ of success and $q = 1-p$, then
\begin{align}
	\pr{X=k} &= \comb{n}{k}p^kq^{n-k}, \quad k = 0,1,\dots,n \label{eq:EE26BTECH11205-binom-pmf}
	\\
	\mean{X} &= \sum_{k=0}^{n}k\,\pr{X=k} = np \label{eq:EE26BTECH11205-binom-mean}
\end{align}
\solution Since $k\comb{n}{k} = n\comb{n-1}{k-1}$,
\begin{align}
	\mean{X} = np\sum_{k=1}^{n}\comb{n-1}{k-1}p^{k-1}q^{n-k} = np\brak{p+q}^{n-1} = np
\end{align}
\item If $\vec{A}$ and $\vec{B}$ are invertible, then
\begin{align}
	\brak{\vec{A}\vec{B}}^{-1} = \vec{B}^{-1}\vec{A}^{-1} \label{eq:EE26BTECH11205-inv-prod}
\end{align}
Also, $\vec{A}^{-1} = \frac{\operatorname{adj}\brak{\vec{A}}}{\mydet{\vec{A}}}$, where $\operatorname{adj}\brak{\vec{A}}$ is the transpose of the matrix of cofactors. \label{eq:EE26BTECH11205-adjugate}

\smallskip
\solution $\brak{\vec{A}\vec{B}}\brak{\vec{B}^{-1}\vec{A}^{-1}} = \vec{A}\brak{\vec{B}\vec{B}^{-1}}\vec{A}^{-1} = \vec{I}$.
\item If $V^{\prime}\brak{h_0} = 0$ and $V^{\prime\prime}\brak{h_0} < 0$, then $V$ has a local maximum at $h_0$. \label{eq:EE26BTECH11205-second-deriv}
\item A right circular cone of height $h$ and base radius $r$ is inscribed in a sphere of radius $R$, with the apex and the base circle on the sphere. The axial section shows that the centre of the base is at distance $\abs{h-R}$ from the centre of the sphere, so
\begin{align}
	r^2 = R^2 - \brak{h-R}^2 = h\brak{2R-h},\; V = \frac{\pi}{3}r^2h = \frac{\pi}{3}h^2\brak{2R-h},\; 0<h<2R \label{eq:EE26BTECH11205-cone-vol}
\end{align}
\item The area under $y = f(x) \ge 0$ between $x=a$ and $x=b$ is
\begin{align}
	\int_{a}^{b}f(x)\,dx \label{eq:EE26BTECH11205-area-int}
\end{align}
For a circle of radius $r$, $\int\sqrt{r^2-x^2}\,dx = \frac{x}{2}\sqrt{r^2-x^2} + \frac{r^2}{2}\sin^{-1}\frac{x}{r} + c$.
\item A definite integral is the limit of a sum,
\begin{align}
	\int_{a}^{b}f(x)\,dx = \lim_{n\to\infty}h\sum_{i=1}^{n}f\brak{a+ih},\quad h = \frac{b-a}{n} \label{eq:EE26BTECH11205-riemann}
\end{align}
\begin{align}
	\sum_{i=1}^{n}i = \frac{n\brak{n+1}}{2},\quad \sum_{i=1}^{n}i^2 = \frac{n\brak{n+1}\brak{2n+1}}{6} \label{eq:EE26BTECH11205-sums}
\end{align}
\item For $t = \tan x$,
\begin{align}
	dt = \brak{1+t^2}dx,\quad \frac{\cos^2x}{\cos^2x + k\sin^2x} = \frac{1}{1+kt^2} \label{eq:EE26BTECH11205-tan-sub}
\end{align}
\item The plane through the points $\vec{P}_1$, $\vec{P}_2$ and perpendicular to the plane $\vec{n}_0^{\top}\vec{x} = c_0$ has the normal $\vec{n}$ with $\vec{n}^{\top}\brak{\vec{P}_2-\vec{P}_1} = 0$ and $\vec{n}^{\top}\vec{n}_0 = 0$. Hence
\begin{align}
	\vec{n} = \brak{\vec{P}_2-\vec{P}_1}\times\vec{n}_0,\quad \vec{n}^{\top}\vec{x} = \vec{n}^{\top}\vec{P}_1 \label{eq:EE26BTECH11205-plane-pts}
\end{align}
The plane contains the line $\vec{x} = \vec{A}_0 + \lambda\vec{d}$ if $\vec{n}^{\top}\vec{A}_0 = \vec{n}^{\top}\vec{P}_1$ and $\vec{n}^{\top}\vec{d} = 0$.
\item (Corner point theorem) In a linear programming problem, the feasible region is convex, and the optimum value of the linear objective $Z = \vec{c}^{\top}\vec{x}$ is attained at a corner point. If the region is unbounded and $Z \ge 0$ on it (for example $\vec{c} \ge \vec{0}$ with $\vec{x}\ge\vec{0}$), the minimum exists and is attained at a corner point. \label{eq:EE26BTECH11205-lpp}
\item (Bayes' theorem) For a partition $\cbrak{A_1,\dots,A_n}$ of the sample space and an event $E$,
\begin{align}
	\pr{A_i|E} = \frac{\pr{A_i}\pr{E|A_i}}{\sum_{j=1}^{n}\pr{A_j}\pr{E|A_j}} \label{eq:EE26BTECH11205-bayes}
\end{align}
\end{enumerate}
